% =====================================================================
% RESULTS-SECTION REWRITE -- paste over the current \section{Results}
% (and its inline Tables 1-7) in the main document.
%
% BEFORE PASTING:
%  1. Preamble needs: \usepackage{booktabs,multirow,placeins}
%     (placeins is new here -- for \FloatBarrier below.)
%  2. Copy the four figure files into the same dir as the main .tex:
%       results/figures/fig_tradeoff.pdf
%       results/figures/fig_backbones.pdf
%       results/figures/fig_ablation.pdf
%       results/figures/fig_robustness.pdf
%     (generated by scripts/fig_experiments.py; re-run that script and
%     these four paths regenerate identically if the numbers change.)
%  3. Move appendix_tables.tex's content in after \appendix (separate file).
%  4. Elsewhere in Methodology, \ref{fig:overview} does not resolve to
%     any existing \label -- only fig:method and fig:population exist.
%     Not fixed here since that section isn't in this file; find/replace
%     fig:overview -> fig:method (or whichever is intended) before compiling.
%
% Two factual fixes applied vs. the earlier draft of this section (both
% checked against results/tables/table2_chain.csv and paper_tables.tex):
%  - The HHAR "head-tuning beats PACT" exception is a Table 2B (frozen-
%    MOMENT chain) fact -- Figure~\ref{fig:tradeoff} plots ONLY the
%    TSLANet chain (Table 2A), where PACT strictly dominates every
%    baseline on every dataset, so that exception is not visible in it.
%    Re-pointed the citation to Appendix Table~\ref{tab:chain-moment}.
%  - "order of magnitude more forgetting" overstated the gap: LwF(head)
%    vs PACT on HHAR/MOMENT is F=0.050 vs 0.017, i.e. ~3x, not ~10x.
%    Corrected in both this file and the Table~\ref{tab:chain-moment}
%    caption in appendix_tables.tex.
%  - "Accuracy scales smoothly with the calibration budget on every
%    dataset" overstated Table 7: Sleep-EDF dips 0.630->0.610 from
%    1%->2% before rising. Softened to name the exception explicitly
%    (this is the same issue flagged for the intro's population numbers
%    in the cleanup list -- worth checking that fix landed there too).
% =====================================================================

\section{Experiments}
\label{sec:experiments}

\subsection{Setup}
\label{sec:setup}
We evaluate on three time-series classification benchmarks spanning
distinct shift types: \textbf{FD} (bearing fault diagnosis; 4 operating
conditions), \textbf{HHAR} (human activity recognition; 9 device/subject
domains), and \textbf{Sleep-EDF} (sleep-stage classification; 17 subject
domains). We use two protocols. The \emph{pairwise} protocol follows
AdaTime: a model is adapted from a single source to a single target
operating condition, averaged over 10 source$\rightarrow$target shifts per
dataset; this is the native setting of the UDA/TTA baselines. The
\emph{domain-incremental chain} protocol---the setting PACT
targets---presents the domains of a dataset sequentially, so that
forgetting on earlier domains is measured directly after the full stream.

We instantiate PACT on two specialist encoders (TSLANet, PatchTST) and
five \emph{frozen} time-series foundation models (MOMENT-1-large,
Chronos-Bolt, Sundial, Time-MoE, Moirai). Baselines fall into four families: unsupervised
domain adaptation (DANN, Deep~CORAL, MMDA, and the semi-supervised
DANN\textsuperscript{\dag}/CDAN\textsuperscript{\dag}/CoTMix\textsuperscript{\dag}),
test-time adaptation (TENT), fine-tuning references (full and head-only),
and continual learning (EWC, MAS, LwF, source-distillation). Every
adapting method receives the same 5\% labeled calibration budget per
domain as PACT. We report AA, BWT, $F$, and RF as defined in
\S\ref{sec:evaluation-metrics} (all on a macro-$F_1$ basis; for $F$ and
RF, $0$ denotes no forgetting). Full per-method numbers for every figure
below are given in Appendix~\ref{app:full-tables}.

\subsection{PACT Retains the Source and Recovers the Drift}
\label{sec:main-results}
Figure~\ref{fig:tradeoff} places every method on the accuracy--forgetting
plane for the domain-incremental chain (TSLANet backbone). The desirable
region is the top-left: high average accuracy with near-zero forgetting.
PACT is the only method that occupies it on all three benchmarks.
Host-modifying baselines are pushed down-and-right along a visible
trade-off: full fine-tuning and LwF recover some accuracy but forget
heavily ($F\!\ge\!0.25$ on FD), whereas the least-forgetting baselines
(EWC, head-tuning) avoid forgetting largely by staying close to the
non-adapted model and never recover the drift. On the pairwise AdaTime
benchmark (Appendix Table~\ref{tab:main-pairwise}), the same pattern holds
against UDA and TTA: PACT attains the best AA on FD and HHAR while holding
$F\!\approx\!0$, matching or exceeding the strongest semi-supervised
baselines that additionally consume target labels. We note one honest
exception, omitted from Figure~\ref{fig:tradeoff} because it is specific
to the frozen-TSFM chain rather than TSLANet: on HHAR with frozen MOMENT
(Appendix Table~\ref{tab:chain-moment}), LwF restricted to the head reaches
marginally higher accuracy than PACT (0.646 vs.\ 0.627) but at roughly
$3\times$ the forgetting (F: 0.050 vs.\ 0.017) -- the one setting in the
paper where a baseline is not strictly dominated on both axes.

\begin{figure*}[t]
\centering
\includegraphics[width=\textwidth]{fig_tradeoff.pdf}
\caption{\textbf{Accuracy--forgetting trade-off on the domain-incremental
chain (TSLANet backbone).} Top-left is best. PACT ($\star$) is the only
method that simultaneously retains source/earlier-domain accuracy and
avoids forgetting; every host-modifying baseline trades one for the other.
Per-method values in Appendix Table~\ref{tab:chain-tslanet}.}
\label{fig:tradeoff}
\end{figure*}
\FloatBarrier

\subsection{PACT Is Backbone-Agnostic, Including Frozen TSFMs}
\label{sec:backbones}
Because PACT only requires an embedding and a frozen head, it applies
unchanged to specialist encoders and to forecasting-pretrained TSFMs used
purely as frozen feature extractors. Figure~\ref{fig:backbones} shows that
near-zero forgetting is preserved across all seven backbones
($F\!\le\!0.022$ everywhere). The two backbone families are complementary:
the frozen TSFMs (MOMENT, Chronos-Bolt, Sundial, Moirai) cluster at the top
on FD (up to $0.984$), while the lighter specialist TSLANet is strongest on
HHAR and Sleep-EDF; Time-MoE trails the other TSFMs on FD ($0.943$) but is
still comfortably above the specialist PatchTST. This confirms that PACT's
guarantee is a property of the architecture, not of any single encoder, and
that a frozen TSFM---which cannot be fine-tuned at all---becomes a
competitive continual classifier under PACT.

\begin{figure*}[t]
\centering
\includegraphics[width=\textwidth]{fig_backbones.pdf}
\caption{\textbf{PACT across seven backbones on the chain.} Top: average
accuracy (AA). Bottom: forgetting $F$, which stays negligible for every
backbone and dataset. Full numbers in Appendix Table~\ref{tab:backbones}.}
\label{fig:backbones}
\end{figure*}
\FloatBarrier

\subsection{Ablation: The Gate Is a Necessity, Not a Refinement}
\label{sec:ablation}
Figure~\ref{fig:ablation}(a) isolates each component on the FD chain.
Accuracy rises monotonically from the frozen head, to full fine-tuning, to
ungated adapters, to full PACT---but only PACT keeps forgetting near zero;
full fine-tuning and ungated adapters recover accuracy at the cost of
large $F$ (up to $0.28$ and $0.11$ respectively, worst-case across
backbones). Adapters supply plasticity; the gate is what converts that
plasticity into forgetting-free accumulation. Figure~\ref{fig:ablation}(b)
explains why. Soft-mixture routing matches or beats hard $\arg\max$
routing and sits just below the oracle (perfect-routing) ceiling. By
Proposition~\ref{prop:frozen} the gap between soft and oracle is the
\emph{sole} channel through which PACT can forget; that gap is
$\le\!0.01$ on FD and Sleep-EDF and widens only on HHAR, precisely where
routing is hardest---localizing the residual forgetting to a single,
interpretable cause rather than to backbone drift.

\begin{figure*}[t]
\centering
\includegraphics[width=\textwidth]{fig_ablation.pdf}
\caption{\textbf{Ablations.} (a)~Component ablation on the FD chain:
accuracy climbs across variants, but only PACT holds $F$ near zero
(annotated). (b)~Gate routing: soft mixture tracks the oracle ceiling; the
residual soft--oracle gap is the only source of forgetting, and is largest
on HHAR. Full numbers in Appendix Tables~\ref{tab:components}
and~\ref{tab:gate}.}
\label{fig:ablation}
\end{figure*}
\FloatBarrier

\subsection{Robustness and Label Efficiency}
\label{sec:robustness}
Figure~\ref{fig:robustness} sweeps PACT's main design choices on the FD
chain. Accuracy increases with the calibration budget on every dataset
(a) -- monotonically on FD and HHAR, and on Sleep-EDF as well apart from a
small $0.630\!\to\!0.610$ dip from 1\% to 2\% (Appendix
Table~\ref{tab:budget-moment}) -- and is stable across concept count $K$
(b) and gate temperature $\tau$ (c); identity vs.\ LoRA adapters and four
domain orderings move AA by at most $0.011$ (Appendix
Table~\ref{tab:hparams}). Forgetting remains $\le\!0.004$ throughout every
sweep. PACT therefore reaches its reported operating point without
careful tuning, and improves predictably when more calibration labels are
available. We also observe the limit of the method transparently: when
calibration on one shift is applied to a structurally different unseen
shift, routing degrades gracefully rather than catastrophically, but the
recovery is partial (Appendix~\ref{app:full-tables}).

\begin{figure*}[t]
\centering
\includegraphics[width=\textwidth]{fig_robustness.pdf}
\caption{\textbf{Robustness and label efficiency (FD chain).}
(a)~Accuracy grows with the calibration budget across datasets and
backbones. (b,c)~Accuracy is flat across concept count $K$ and gate
temperature $\tau$. Forgetting stays $\le\!0.004$ in all runs. Full
numbers in Appendix Tables~\ref{tab:hparams} and~\ref{tab:budget-moment}.}
\label{fig:robustness}
\end{figure*}
\FloatBarrier
